\documentclass[dvipdfmx]{jsarticle}
\usepackage[a4paper, top=20truemm, bottom=20truemm, left=20truemm, right=20truemm]{geometry}
\usepackage{amsmath, amssymb, amsfonts, mathtools, amsthm}
\usepackage[dvipdfmx]{graphicx}
\usepackage[dvipdfmx, unicode=true, colorlinks=true, linkcolor=blue, citecolor=blue]{hyperref}

\numberwithin{equation}{section}

\theoremstyle{plain}
\newtheorem{thm}{定理}[section]
\newtheorem{lem}[thm]{補題}
\newtheorem{ex}[thm]{例}

\makeatletter
\renewenvironment{proof}[1][\proofname]{\par
  \pushQED{\qed}%
  \normalfont \topsep6\p@\@plus6\p@\relax
  \trivlist
  \item\relax
  {\bfseries
  #1\@addpunct{.}}\hspace\labelsep\ignorespaces
}{%
  \popQED\endtrivlist\@endpefalse
}
\makeatother
\renewcommand{\proofname}{証明}

\title{\bf 大数の弱法則の構成}
\author{}
\date{}

\begin{document}
\maketitle

\section{無相関な確率変数列と分散}

\begin{thm}[定理2.2.1]
$E(X_i^2) < \infty$ である確率変数列 $X_1,\dots,X_n$ が無相関であれば，
\[
  \mathrm{var}(X_1 + \cdots + X_n) = \mathrm{var}(X_1) + \cdots + \mathrm{var}(X_n)
\]
が成り立つ．
\end{thm}

\begin{proof}
$\mu_i = E(X_i)$，$S_n = \sum_{i=1}^n X_i$ とおく．$E(S_n) = \sum_{i=1}^n \mu_i$ だから，分散の定義を用い，和の2乗を積の形で書いて展開すると，
\begin{align*}
  \mathrm{var}(S_n)
  &= E\!\left((S_n - E(S_n))^2\right)
   = E\!\left(\left\{\sum_{i=1}^n (X_i - \mu_i)\right\}^2\right) \\
  &= E\!\left(\sum_{i=1}^n \sum_{j=1}^n (X_i - \mu_i)(X_j - \mu_j)\right) \\
  &= \sum_{i=1}^n E\!\left((X_i - \mu_i)^2\right) + 2\sum_{i=1}^n \sum_{j=1}^{i-1} E\!\left((X_i - \mu_i)(X_j - \mu_j)\right)
\end{align*}
となる．ただし，最後の等式では対角の項 $i=j$ を分け，$1 \le i < j \le n$ に関する和と $1 \le j < i \le n$ に関する和が等しいことを用いた．

第1項は $\mathrm{var}(X_1) + \cdots + \mathrm{var}(X_n)$ だから，第2項が $0$ であることを示せばよい．これには $X_i$ と $X_j$ が無相関であることから，次に注意すればよい．
\[
  E\!\left((X_i - \mu_i)(X_j - \mu_j)\right)
  = E(X_iX_j) - \mu_i E(X_j) - \mu_j E(X_i) + \mu_i\mu_j
  = E(X_iX_j) - \mu_i\mu_j = 0
\]
\end{proof}

\begin{lem}[補題2.2.2]
$p > 0$ に対して $E(|Z_n|^p) \to 0$ ならば，$Z_n$ は $0$ に確率収束する．
\end{lem}

\begin{proof}
Chebyshev の不等式を $\varphi(x) = x^p$，$X = |Z_n|$ として用いると，$\epsilon > 0$ に対して
\[
  P(|Z_n| \ge \epsilon) \le \epsilon^{-p} E(|Z_n|^p) \to 0
\]
がわかる．
\end{proof}

\texorpdfstring{$L^2$}{L2}
\begin{thm}[定理2.2.3（$L^2$弱法則）]
$X_1, X_2, \dots$ を，$E(X_i) = \mu$，$\mathrm{var}(X_i) \le C < \infty$ をみたす無相関な確率変数列とする．このとき $S_n = X_1 + \cdots + X_n$ とおくと，$n \to \infty$ のとき $S_n/n$ は $\mu$ に2次平均収束かつ確率収束する．
\end{thm}

\begin{proof}
2次平均収束は，$E(S_n/n) = \mu$，および
\[
  E\!\left(\left(\frac{S_n}{n} - \mu\right)^2\right)
  = \mathrm{var}\!\left(\frac{S_n}{n}\right)
  = \frac{1}{n^2}\left(\mathrm{var}(X_1) + \cdots + \mathrm{var}(X_n)\right)
  \le \frac{Cn}{n^2} \to 0
\]
からわかる．確率収束は，補題2.2.2 を $Z_n = S_n/n - \mu$ に適用すれば得られる．
\end{proof}

定理2.2.3 の特別な場合で最も重要なものは，$X_1, X_2, \dots$ が同じ分布をもつ独立確率変数列，すなわち i.i.d.\ の場合である．この場合，定理2.2.3 は，$E(X_i^2) < \infty$ であれば $S_n/n$ が $n \to \infty$ のとき $\mu = E(X_i)$ に確率収束することを示している．後の定理2.2.14 で，このことが成り立つには $E(|X_i|) < \infty$ であれば十分であることを示す．

\section{応用例：多項式近似（Bernstein 多項式）}

\begin{ex}[例2.2.4（多項式近似）]
$f$ を $[0,1]$ 上の連続関数とし，$f_n$ を
\[
  f_n(x) = \sum_{m=0}^n \binom{n}{m} x^m (1-x)^{n-m} f\!\left(\frac{m}{n}\right), \qquad
  \binom{n}{m} = \frac{n!}{m!(n-m)!}
\]
で与えられる $f$ に対応する $n$ 次 Bernstein 多項式とすると，$n \to \infty$ のとき次が成り立つ．
\[
  \sup_{x \in [0,1]} |f_n(x) - f(x)| \to 0
\]
\end{ex}

\begin{proof}
まず，$S_n$ を $P(X_i = 1) = p$，$P(X_i = 0) = 1-p$ である独立確率変数の和とすると，$E(X_i) = p$，$\mathrm{var}(X_i) = p(1-p)$ であり，
\[
  P(S_n = m) = \binom{n}{m} p^m (1-p)^{n-m}
\]
が成り立つので，$E(f(S_n/n)) = f_n(p)$ である．定理2.2.3 より，$n \to \infty$ のとき $S_n/n$ は $p$ に確率収束する．これらが $f_n(p)$ を考える動機であるが，この証明では，定理2.2.3 の結果以上にその証明を用いる．

$p \in [0,1]$ に対して $p(1-p) \le 1/4$ であることを用いると，定理2.2.3 の証明から
\[
  P\!\left(\left|\frac{S_n}{n} - p\right| > \delta\right)
  \le \frac{\mathrm{var}(S_n/n)}{\delta^2}
  = \frac{p(1-p)}{n\delta^2}
  \le \frac{1}{4n\delta^2}
\]
を得る．$E(f(S_n/n)) \to f(p)$ の証明には，$M = \sup_{x \in [0,1]} |f(x)|$ とおき，$\epsilon > 0$ に対して $\delta > 0$ を，$|x - y| < \delta$ ならば $|f(x) - f(y)| < \epsilon$ となるようにとる（連続関数は有界区間上で一様連続だから可能である）．ここで Jensen の不等式を用いると
\[
  \left|E\!\left(f\!\left(\frac{S_n}{n}\right)\right) - f(p)\right|
  \le E\!\left(\left|f\!\left(\frac{S_n}{n}\right) - f(p)\right|\right)
  \le \epsilon + 2M P\!\left(\left|\frac{S_n}{n} - p\right| > \delta\right)
\]
がわかる．よって，$n \to \infty$ とすると，$\limsup_{n\to\infty} |E(f(S_n/n)) - f(p)| \le \epsilon$ となり，$\epsilon$ は任意なので結論を得る．
\end{proof}

\section{三角配列に対する弱法則}

確率論における多くの古典的な結果は独立同分布（i.i.d.）の場合に扱われるが，行に関する和 $S_n = X_{n,1} + \cdots + X_{n,n}$ を考える場合，各行の確率変数は独立でありさえすれば，同分布であるという仮定なしに成り立つ結果も多い．このような配列を三角配列（triangular array）とよぶ．

確率変数の弱法則を2次のモーメントが存在しない場合に拡張するためには，打切りをして Chebyshev の不等式を用いる．まず，一般的で役に立つ定理を述べる．証明に必要なことをそのまま仮定するので，証明は容易であるが，後述の特別な場合に仮定を確認するためには少し議論が必要である．一般的な結果を述べることで，証明の本質的な部分を特定できる．

確率変数 $X$ の打切り（truncation）を
\[
  \bar{X} = X 1_{(|X| \le M)} =
  \begin{cases}
    X & (|X| \le M) \\
    0 & (|X| > M)
  \end{cases}
\]
とする．

\begin{thm}[定理2.2.11（三角配列に対する弱法則）]
各 $n$ に対し，$X_{n,k}$，$1 \le k \le n$ は独立とする．$b_n > 0$，$b_n \to \infty$ をみたす $b_n$ に対して $\bar{X}_{n,k} = X_{n,k} 1_{(|X_{n,k}| \le b_n)}$ とおき，$n \to \infty$ のとき次が成り立つと仮定する．
\begin{enumerate}
  \item[(i)] $\displaystyle \sum_{k=1}^n P(|X_{n,k}| > b_n) \to 0$
  \item[(ii)] $\displaystyle b_n^{-2} \sum_{k=1}^n E(\bar{X}_{n,k}^2) \to 0$
\end{enumerate}
このとき，$S_n = X_{n,1} + \cdots + X_{n,n}$ とし，$a_n = \sum_{k=1}^n E(\bar{X}_{n,k})$ とおくと，
\[
  \frac{S_n - a_n}{b_n} \to 0
\]
が確率収束の意味で成り立つ．
\end{thm}

\begin{proof}
$\bar{S}_n = \bar{X}_{n,1} + \cdots + \bar{X}_{n,n}$ とおく．明らかに
\[
  P\!\left(\left|\frac{S_n - a_n}{b_n}\right| > \epsilon\right)
  \le P(S_n \ne \bar{S}_n) + P\!\left(\left|\frac{\bar{S}_n - a_n}{b_n}\right| > \epsilon\right)
\]
が成り立つ．第1項については，仮定(i)より
\[
  P(S_n \ne \bar{S}_n)
  \le P\!\left(\bigcup_{k=1}^n \{\bar{X}_{n,k} \ne X_{n,k}\}\right)
  \le \sum_{k=1}^n P(|X_{n,k}| > b_n) \to 0
\]
が成り立つ．第2項については，Chebyshev の不等式，$a_n = E(\bar{S}_n)$ であること，定理2.2.1 および $\mathrm{var}(X) \le E(X^2)$ より，(ii)を用いると
\[
  P\!\left(\left|\frac{\bar{S}_n - a_n}{b_n}\right| > \epsilon\right)
  \le \frac{1}{\epsilon^2} E\!\left(\left(\frac{\bar{S}_n - a_n}{b_n}\right)^2\right)
  = \frac{1}{\epsilon^2 b_n^2}\, \mathrm{var}(\bar{S}_n)
  = \frac{1}{\epsilon^2 b_n^2} \sum_{k=1}^n \mathrm{var}(\bar{X}_{n,k})
  \le \frac{1}{\epsilon^2 b_n^2} \sum_{k=1}^n E(\bar{X}_{n,k}^2) \to 0
\]
となり，定理の結論を得る．
\end{proof}

\section{大数の弱法則}

\begin{thm}[定理2.2.12（大数の弱法則）]
$X_1, X_2, \dots$ を
\[
  xP(|X_i| \ge x) \to 0, \qquad x \to \infty
\]
をみたす i.i.d.\ とする．このとき，$S_n = X_1 + \cdots + X_n$，$\mu_n = E(X_1 1_{(|X_1| \le n)})$ とおくと，$S_n/n - \mu_n$ は $0$ に確率収束する．
\end{thm}

\emph{注意．} 定理の仮定は，$S_n/n - a_n \to 0$ となる定数 $a_n$ の存在のために必要である．

\begin{proof}
定理2.2.11 を $X_{n,k} = X_k$，$b_n = n$ として適用する．条件(i)を確認するために，仮定から
\[
  \sum_{k=1}^n P(|X_{n,k}| > n) = n P(|X_i| > n) \to 0
\]
が成り立つことに注意する．(ii)の確認には，$n^{-2} \cdot n E(\bar{X}_{n,1}^2) \to 0$ を示す必要がある．このために，今後繰り返し用いる次の結果を示す．
\end{proof}

\begin{lem}[補題2.2.13]
$Y \ge 0$，$p > 0$ ならば，
\[
  E(Y^p) = \int_0^\infty p y^{p-1} P(Y > y)\, dy
\]
が成り立つ．
\end{lem}

\begin{proof}
期待値の定義と（非負な確率変数に対する）Fubini の定理を用いて積分を計算すると，次を得る．
\begin{align*}
  \int_0^\infty p y^{p-1} P(Y > y)\, dy
  &= \int_0^\infty \int_\Omega p y^{p-1} 1_{(Y>y)}\, dP\, dy
   = \int_\Omega \int_0^\infty p y^{p-1} 1_{(Y>y)}\, dy\, dP \\
  &= \int_\Omega \int_0^Y p y^{p-1}\, dy\, dP
   = \int_\Omega Y^p\, dP = E(Y^p)
\end{align*}
\end{proof}

定理2.2.12 の証明に戻る．補題2.2.13 と $\bar{X}_{n,1} = X_1 1_{(|X_1| \le n)}$ より，$P(|\bar{X}_{n,1}| > y)$ は $y \ge n$ ならば $0$ で，$y \le n$ ならば $P(|X_1| > y) - P(|X_1| > n)$ に等しく，
\[
  E(\bar{X}_{n,1}^2) = \int_0^\infty 2y P(|\bar{X}_{n,1}| > y)\, dy \le \int_0^n 2y P(|X_1| > y)\, dy
\]
が成り立つことに注意する．$yP(|X_1| > y) \to 0$ を用いて，$n \to \infty$ のとき
\[
  \frac{E(\bar{X}_{n,1}^2)}{n} = \frac{1}{n} \int_0^n 2y P(|X_1| > y)\, dy \to 0
\]
が成り立つことを示したい．右辺は $g(y) = 2yP(|X_1| > y)$ の $[0,n]$ における平均であり，$y \to \infty$ のとき $g(y) \to 0$ だから，この収束は直感的に理解できる．詳しく述べるために，$0 \le g(y) \le 2y$ かつ $y \to \infty$ のとき $g(y) \to 0$ だから $\sup g(y) < \infty$ であることに注意して，$g_n(y) = g(ny)$ とおく．$g_n$ は有界で $g_n(y) \to 0$ a.e.\ だから，有界収束定理より
\[
  \frac{1}{n} \int_0^n g(y)\, dy = \int_0^1 g_n(x)\, dx \to 0
\]
となって，定理の結論を得る．

\emph{注意．} 補題2.2.13 を $p = 1 - \epsilon$，$\epsilon > 0$ として適用すると，$xP(|X_1| > x) \to 0$ より $E(|X_1|^{1-\epsilon}) < \infty$ となるので，定理2.2.12 の仮定が平均の有限性よりずっと弱いわけではない．

最後に，最もよく知られた形の弱法則を述べる．

\begin{thm}[定理2.2.14]
$X_1, X_2, \dots$ を $E(|X_i|) < \infty$ である i.i.d.\ とし，$S_n = X_1 + \cdots + X_n$，$\mu = E(X_1)$ とおく．このとき，$S_n/n$ は $\mu$ に確率収束する．
\end{thm}

\begin{proof}
優収束定理から
\[
  xP(|X_1| > x) \le E\!\left(|X_1| 1_{(|X_1|>x)}\right) \to 0, \qquad x \to \infty
\]
\[
  \mu_n = E\!\left(X_1 1_{(|X_1| \le n)}\right) \to E(X_1) = \mu, \qquad n \to \infty
\]
がわかる．定理2.2.12 を用いると，$\epsilon > 0$ に対して $P(|S_n/n - \mu_n| > \epsilon/2) \to 0$ もわかる．$\mu_n \to \mu$ だから，$P(|S_n/n - \mu| > \epsilon) \to 0$ を得る．
\end{proof}

\section{弱法則が成り立たない例：Cauchy 分布}

\begin{ex}[例2.2.15]
弱法則が成り立たない例を挙げるために，$X_1, X_2, \dots$ が独立で次の Cauchy 分布に従うとする．
\[
  P(X_i \le x) = \int_{-\infty}^x \frac{dt}{\pi(1+t^2)}
\]
$x \to \infty$ のとき
\[
  P(|X_1| > x) = 2\int_x^\infty \frac{dt}{\pi(1+t^2)} \sim \frac{2}{\pi}\int_x^\infty t^{-2}\, dt = \frac{2}{\pi} x^{-1}
\]
である．このことから，$S_n/n - \mu_n \to 0$ をみたす数列 $\mu_n$ が存在しないと結論づけることができる．実際，$S_n/n$ が $X_1$ と同じ分布をもつことが示される．
\end{ex}

次の例が示すように，$E(|X_i|) = \infty$ のときにも弱法則が成り立つ場合がある．

\section{応用例：サンクトペテルブルクのパラドックス}

\begin{ex}[例2.2.16（サンクトペテルブルクのパラドックス）]
$X_1, X_2, \dots$ を，$j \ge 1$ に対し
\[
  P(X_i = 2^j) = 2^{-j}
\]
である独立確率変数とする．言い換えると，硬貨投げで $j$ 回続けて表が出ると $2^j$ ドル獲得する場合である．このパラドックスは，期待値は $E(X_1) = \infty$ だが，このゲームに無限の参加費は払わないだろうということである．定理2.2.11 を用いると，ゲーム（硬貨投げ）を $n$ 回行うとしたときの参加費がわかる．
\end{ex}

\begin{proof}[考察]
この例では，$X_{n,k} = X_k$ とおく．定理2.2.11 を適用するには $b_n$ を選ぶ必要がある．そのためには，条件(ii)を確認する際には $b_n$ をできるだけ小さくとり，しかも(i)をみたさなければならない．このことを念頭に，整数 $m$ に対して
\[
  P(X_1 \ge 2^m) = \sum_{j=m}^\infty 2^{-j} = 2^{-m+1}
\]
であることに注意する．$m(n) = \log_2 n + K(n)$ とおく．ただし，$K(n)$ は $K(n) \to \infty$ かつ $m(n)$ が整数になる（よって，上が成り立つ）ように選ぶ．$b_n = 2^{m(n)}$ とおくと，
\[
  nP(X_1 \ge b_n) = n 2^{-m(n)+1} = 2^{-K(n)+1} \to 0
\]
となり，(i)が成り立つ．(ii)を確認するために，$\bar{X}_{n,k} = X_k 1_{(|X_k| \le b_n)}$ とおくと
\[
  E(\bar{X}_{n,k}^2) = \sum_{j=1}^{m(n)} 2^{2j} \cdot 2^{-j} \le 2^{m(n)} \sum_{k=0}^\infty 2^{-k} = 2b_n
\]
であることに注意する．(ii)で考える量は $2n/b_n$ より小さく，
\[
  b_n = 2^{m(n)} = n 2^{K(n)} \qquad \text{かつ} \qquad K(n) \to \infty
\]
だから $0$ に収束する．

次の2段階では，$a_n$ を求めて，定理2.2.11 を適用する．
\[
  E(\bar{X}_{n,k}) = \sum_{j=1}^{m(n)} 2^j 2^{-j} = m(n)
\]
より，$a_n = n m(n)$ である．$m(n) = \log n + K(n)$ だから（この例の中ではずっと対数の底は $2$ である），$n \to \infty$ のとき $K(n)/\log n \to 0$ ならば $a_n/(n\log n) \to 1$ が成り立つ．そして，定理2.2.11 を適用すると，確率収束の意味で
\[
  \frac{S_n - a_n}{n 2^{K(n)}} \to 0
\]
を得る．$n$ が大きいとき $K(n) \le \log\log n$ であるようにとると，この収束が分母を $n\log n$ に置き換えても成立し，$S_n/(n\log n)$ が $1$ に確率収束する．
\end{proof}

\end{document}